\section{Oversight: am I still in control?}
\label{sec:oversight}

The clinician must remain the responsible agent. The fourth question asks whether
the human stays in control: how to override, pause, or stop the system, and when it
escalates a decision to a qualified person rather than acting alone.

[DRAFTING INSTRUCTIONS: in the full-framework, describe human-in-the-loop and
human-on-the-loop supervision, the ESCALATE path, and the unconditional override,
grounded in \texttt{lee2004trust} (appropriate reliance) and
\texttt{goddard2012automation} (automation bias in clinical decision support).
Reproduce Mermaid
\texttt{adoption/mermaid/diagrams/09-oversight-loop-state.md},
\texttt{14-escalation-pathway-sequence.md}, and \texttt{19-override-stop-state.md}
as colored cfwfig TikZ. Carry the KEY TERMS ``human-in-the-loop'' and
``human-on-the-loop''. Build a body-width table of oversight modes and the
clinician's controls.]

\begin{cfwfig}
\node[cftwo] (a) at (0,0) {Monitor}; \node[cfact] (b) at (3.2,0) {Recommend}; \node[cfdec] (c) at (6.6,0) {Decision};
\node[cfhope] (d) at (9.8,1.0) {Execute}; \node[cfthree] (e) at (9.8,-0.2) {Escalate}; \node[cfrisk] (f) at (9.8,-1.4) {Override};
\draw[cfedge] (a)--(b); \draw[cfedge] (b)--(c);
\draw[cfedge] (c)--(d); \draw[cfedge] (c)--(e); \draw[cfedge] (c)--(f);
\end{cfwfig}
\vspace{-0.2cm}
\cfwcaption{Figure 5. The oversight loop (placeholder; expanded from Mermaid 09 in
the full-framework).}

Control is not a slogan but a set of switches: a pause, an override, and an
escalation path that the clinician can reach at any moment.
